\newpage
\section*{Solucionario completo}

\noindent Cada respuesta incluye la clave y una justificación breve. En los ejercicios numéricos se indica la ecuación o el cálculo principal.

\subsection*{Bloque A. Fundamentos teóricos y conceptuales}

\sol{1}{C}{La RMS es la pendiente (en valor absoluto) de la curva de indiferencia: indica cuántas unidades de $y$ se está dispuesto a ceder por una unidad adicional de $x$ sin alterar la utilidad. No es la pendiente de la recta presupuestaria ni una relación de precios.}

\sol{2}{A}{$UMg_x=2xy^3$ y $UMg_y=3x^2y^2$. Por tanto, $RMS_{xy}=UMg_x/UMg_y=2y/(3x)$.}

\sol{3}{D}{Una transformación monótona creciente conserva el orden de las canastas, las curvas de indiferencia y la RMS; las utilidades marginales se multiplican por la derivada de la transformación, por lo que pueden cambiar.}

\sol{4}{B}{La convexidad de las preferencias equivale a curvas de indiferencia convexas al origen y a una RMS decreciente a lo largo de la curva.}

\sol{5}{C}{El intercepto vertical $I/p_y$ no cambia, pero el intercepto horizontal $I/p_x$ disminuye: la recta gira hacia adentro sobre el eje $x$.}

\sol{6}{B}{El impuesto eleva el precio efectivo de $x$ de $p_x$ a $p_x+t$; la restricción es $(p_x+t)x+p_yy=I$.}

\sol{7}{D}{El subsidio en especie obliga a consumir al menos $G$ unidades de $x$. Si el consumidor valora poco ese bien, la canasta preferida con dinero en efectivo no es alcanzable; el efectivo es, en general, débilmente preferido.}

\sol{8}{A}{En un óptimo interior se requiere $RMS_{xy}=p_x/p_y$ y que la condición de segundo orden (curvatura) garantice un máximo. La igualdad de utilidades marginales solo sería válida si los precios fueran iguales.}

\sol{9}{C}{La utilidad por sol gastado es mayor en $x$ que en $y$. Reasignar gasto hacia $x$ aumenta la utilidad hasta que se igualan las utilidades marginales por sol.}

\sol{10}{B}{$UMg_x/p_x=a/p_x$ y $UMg_y/p_y=b/p_y$. Si $p_x/p_y<a/b$, entonces $a/p_x>b/p_y$: todo el ingreso se destina a $x$.}

\sol{11}{D}{Los complementos perfectos tienen curvas de indiferencia en forma de L, con vértice en la combinación que satisface la proporción fija.}

\sol{12}{A}{Con $U=\min(x,y)$, el óptimo exige $x=y$. Un mayor $p_x$ reduce la cantidad demandada de ambos bienes en la misma proporción.}

\sol{13}{B}{La condición $v'(x)=p_x/p_y$ determina $x$ sin depender de $I$, siempre que el consumo de $y$ sea positivo. El ingreso se ajusta principalmente vía $y$.}

\sol{14}{C}{Con preferencias homotéticas, la curva de ingreso-consumo es un rayo desde el origen, las elasticidades ingreso son unitarias y las participaciones del gasto son constantes.}

\sol{15}{D}{Un bien inferior reduce su consumo cuando aumenta el ingreso: $\varepsilon_I<0$.}

\sol{16}{A}{Un bien de lujo aumenta su participación en el gasto cuando sube el ingreso: $\varepsilon_I>1$.}

\sol{17}{C}{La demanda marshalliana surge de maximizar $U(x,y)$ sujeto a $p_xx+p_yy=I$.}

\sol{18}{B}{La demanda hicksiana surge de minimizar $p_xx+p_yy$ sujeto a $U(x,y)=u$.}

\sol{19}{D}{$e(p,u)=\min\{p\cdot x:U(x)=u\}$ es el gasto mínimo para alcanzar la utilidad $u$ a los precios $p$.}

\sol{20}{A}{La identidad de Roy establece $x_i(p,I)=-\dfrac{\partial v/\partial p_i}{\partial v/\partial I}$.}

\sol{21}{C}{El lema de Shephard establece que la derivada de la función de gasto respecto del precio es la demanda hicksiana: $\partial e/\partial p_i=h_i(p,u)$.}

\sol{22}{B}{La ecuación de Slutsky separa el efecto total de un cambio de precio en efecto sustitución (movimiento a lo largo de la curva de indiferencia) y efecto ingreso (desplazamiento a otra curva).}

\sol{23}{D}{Para un bien normal, el efecto ingreso refuerza el efecto sustitución; por ello la demanda marshalliana es más elástica (más plana) y la hicksiana es más empinada.}

\sol{24}{A}{El bien Giffen es un caso extremo de bien inferior en el que el efecto ingreso negativo domina al efecto sustitución, generando una relación precio-cantidad positiva.}

\sol{25}{C}{Una demanda vertical implica que la cantidad no responde al precio: $\varepsilon=0$.}

\sol{26}{B}{Si $|\varepsilon|>1$, el aumento porcentual del precio supera la caída porcentual de la cantidad; el ingreso total $PQ$ disminuye.}

\sol{27}{D}{La demanda marshalliana es homogénea de grado cero en $(p,I)$: multiplicar todos los precios y el ingreso por una misma constante no altera la canasta óptima.}

\sol{28}{A}{La ley de Walras establece $\sum_i p_i z_i(p)=0$: el valor agregado de los excesos de demanda es cero. Si todos los mercados menos uno se vacían, el restante también.}

\sol{29}{C}{Las preferencias lexicográficas no satisfacen continuidad y no pueden representarse mediante una función de utilidad continua.}

\sol{30}{B}{WARP: si $x$ se revela preferida a $y$ cuando ambas son asequibles, no puede ocurrir simultáneamente que $y$ se revele preferida a $x$ en otra situación en que ambas sean asequibles.}

\subsection*{Bloque B. Aplicaciones cuantitativas y optimización}

\sol{31}{D}{Con Cobb-Douglas simétrica, $x=\dfrac{0.5I}{p_x}=\dfrac{0.5(200)}{4}=25$.}

\sol{32}{A}{$y=\dfrac{(2/3)I}{p_y}=\dfrac{(2/3)(300)}{10}=20$. El gasto en $y$ es $p_yy=10(20)=200$.}

\sol{33}{C}{$UMg_x/p_x=4/2=2$ y $UMg_y/p_y=3/3=1$. Como $2>1$, el consumidor destina todo el ingreso a $x$: $x=100/2=50$, $y=0$.}

\sol{34}{B}{En el óptimo $y=2x$. La restricción es $3x+2(2x)=7x=140$, de donde $x=20$, $y=40$ y $U=\min(40,40)=40$.}

\sol{35}{D}{En el óptimo $x=2y$. La restricción es $4(2y)+2y=10y=120$, de donde $y=12$ y $x=24$.}

\sol{36}{A}{La función es CES con $\rho=0.5$. La elasticidad de sustitución es $\sigma=\dfrac{1}{1-\rho}=\dfrac{1}{0.5}=2$.}

\sol{37}{C}{Aquí $\rho=-1$. Entonces $\sigma=\dfrac{1}{1-(-1)}=0.5$.}

\sol{38}{B}{$RMS=UMg_x/UMg_y=(1/x)/1=1/x$. Igualando a $p_x/p_y=2$: $x=0.5$.}

\sol{39}{D}{$RMS=UMg_x/UMg_y=1/(1/y)=y$. Igualando a $p_x/p_y=1/2$: $y=0.5$.}

\sol{40}{A}{En Cobb-Douglas, la participación de $x$ es $\dfrac{a}{a+b}=\dfrac{0.3}{0.3+0.6}=\dfrac{1}{3}\approx33.3\%$.}

\sol{41}{C}{En Cobb-Douglas, la demanda de $x$ es $x=\dfrac{aI}{p_x}$, por lo que $\varepsilon_{x,p_x}=-1$.}

\sol{42}{B}{La demanda de $x$ no depende de $p_y$ en Cobb-Douglas; la elasticidad cruzada es cero.}

\sol{43}{D}{El gasto de subsistencia es $p_x\gamma_x+p_y\gamma_y=2(10)+4(5)=40$. El ingreso residual es $100-40=60$. Con participaciones iguales: $x=10+\dfrac{0.5(60)}{2}=25$.}

\sol{44}{A}{$y=5+\dfrac{0.5(60)}{4}=12.5$. El gasto es $4(12.5)=50$.}

\sol{45}{C}{Inicialmente $x=60$. Para que el consumidor pueda comprar la canasta original a los nuevos precios se requiere $I'=2(60)+1(60)=180$. Con $I'=180$ y $p_x=2$, $x=45$. El efecto sustitución es $45-60=-15$.}

\sol{46}{B}{La nueva demanda marshalliana es $x=120/(2\cdot2)=30$. El efecto ingreso es $30-45=-15$. El efecto total es $-30$.}

\sol{47}{D}{La utilidad inicial es $U_0=\sqrt{50\cdot50}=50$. La demanda hicksiana es $x=U_0\sqrt{p_y/p_x}=50\sqrt{1/4}=25$ e $y=U_0\sqrt{p_x/p_y}=100$. El gasto mínimo a los nuevos precios es $4(25)+100=200$. La CV es $200-100=100$.}

\sol{48}{A}{La nueva canasta es $x=12.5$, $y=50$, con $U_1=\sqrt{625}=25$. A los precios originales, el gasto mínimo para alcanzar $U_1$ es $e(1,1,25)=2(25)=50$. La EV es $50-100=-50$; en valor absoluto, la disposición a aceptar es 50.}

\sol{49}{C}{Con $P=60$, $Q=20$; con $P=80$, $Q=10$. La pérdida de excedente es el área del trapecio: $\dfrac{20+10}{2}(80-60)=300$.}

\sol{50}{B}{$UMg_x/UMg_y=1/\sqrt{x}$. Igualando a $1$: $x=1$; $y=99$; $U_0=101$. Con $p_x=4$, la demanda hicksiana para $U_0=101$ satisface $1/\sqrt{x}=4$, es decir $x=1/16$; entonces $y=101-2(1/4)=100.5$. El gasto mínimo es $4(1/16)+100.5=100.75$. La CV es $100.75-100=0.75$.}

\sol{51}{D}{$Q=100-2(20)=60$. Entonces $\varepsilon=-2\left(\dfrac{20}{60}\right)=-0.667\approx-0.67$.}

\sol{52}{A}{Para $Q=AP^{-\alpha}$ la elasticidad es constante e igual a $-\alpha$; aquí $-2$.}

\sol{53}{C}{$\Delta Q\%\approx\varepsilon\cdot\Delta P\%=(-0.5)(10\%)=-5\%$.}

\sol{54}{B}{$\Delta Q\%\approx\varepsilon_I\cdot\Delta I\%=(0.8)(10\%)=8\%$.}

\sol{55}{D}{La elasticidad cruzada positiva indica bienes sustitutos. $\Delta Q\%=(0.3)(10\%)=3\%$.}

\sol{56}{A}{$UMg_x/p_x=2/3\approx0.667$ y $UMg_y/p_y=1/2=0.5$. Como $0.667>0.5$, todo el ingreso se destina a $x$: $x=60/3=20$, $y=0$.}

\sol{57}{C}{En el óptimo $x=3y$. La restricción es $2(3y)+1y=7y=70$, de donde $y=10$ y $x=30$.}

\sol{58}{B}{$x=I/(2p_x)=160/4=40$ e $y=I/(2p_y)=160/16=10$. En el óptimo, $RMS=y/x=10/40=1/4=p_x/p_y$.}

\sol{59}{D}{Con $U=x^2y$, las participaciones son $2/3$ para $x$ y $1/3$ para $y$. Entonces $x=\dfrac{(2/3)(90)}{3}=20$.}

\sol{60}{A}{En Leontief los bienes se consumen en proporciones fijas; no hay sustitución posible, por lo que $\sigma=0$.}

\sol{61}{C}{Con $U=\sqrt{xy}$, $x=I/(2p_x)$ e $y=I/(2p_y)$. Entonces $v=\sqrt{xy}=\sqrt{\dfrac{I^2}{4p_xp_y}}=\dfrac{I}{2\sqrt{p_xp_y}}$.}

\sol{62}{B}{Para $U=\sqrt{xy}$, $e(p,u)=2u\sqrt{p_xp_y}$. Con $u=10$, $p_x=4$ y $p_y=1$: $e=2(10)(2)=40$.}

\sol{63}{D}{$x_0=100/(2\cdot5)=10$; $x_1=200/(2\cdot5)=20$. El aumento es 10.}

\sol{64}{A}{Con $U=xy$, $y=I/(2p_y)$ no depende de $p_x$. Al variar $p_x$, $y$ permanece constante: la curva de precio-consumo es horizontal.}

\sol{65}{C}{$x=I/(2p_x)$ e $y=I/(2p_y)$; por tanto $y/x=p_x/p_y$ es constante. La curva de ingreso-consumo es una recta que pasa por el origen.}

\subsection*{Bloque C. Casos con implicancias tarifarias y regulatorias}

\sol{66}{B}{Factura $=5+10(1.50)+10(2.50)+5(4.00)=5+15+25+20=65$ soles.}

\sol{67}{D}{Tarifa media $=65/25=2.60$ soles por m$^3$. Es distinta de la tarifa marginal del último bloque (S/ 4.00).}

\sol{68}{A}{El m$^3$ adicional se cobra a la tarifa marginal del bloque en que se encuentra el usuario: S/ 4.00.}

\sol{69}{C}{El cargo fijo es una transferencia lump-sum que reduce el ingreso disponible. Con preferencias cuasilineales, la demanda de agua no depende del ingreso (en el tramo interior), por lo que el efecto inmediato sobre $q$ es nulo.}

\sol{70}{B}{En una restricción quebrada, si la RMS en el punto de quiebre está entre los precios marginales de los bloques adyacentes, ninguna tangencia interior es óptima y el consumidor elige exactamente el límite del bloque.}

\sol{71}{D}{$Q=20-0.5(10)=15$. $\varepsilon=-0.5\left(\dfrac{10}{15}\right)=-0.333\approx-0.33$.}

\sol{72}{A}{Ingreso inicial $=10(15)=150$. Con $P=12$, $Q=14$ e ingreso $=12(14)=168$. Aumenta en S/ 18.}

\sol{73}{C}{Con demanda inelástica, la cantidad cae menos que proporcionalmente que el precio; el ingreso tarifario aumenta. El consumo se reduce levemente.}

\sol{74}{B}{La regla Ramsey-Boiteaux fija mayores márgenes sobre el costo marginal a los grupos con demanda menos elástica, porque su consumo responde menos y la distorsión es menor.}

\sol{75}{D}{Un cargo fijo alto puede excluir a usuarios de baja disposición a pagar, aunque permita cubrir costos fijos. Esto plantea un conflicto entre eficiencia y asequibilidad.}

\sol{76}{A}{Con costo medio decreciente, $P=CMe$ cubre costos (beneficio nulo), pero como $CMe>CMg$, el precio supera el costo marginal y la cantidad es menor que la eficiente.}

\sol{77}{C}{Con $P=CMg<CMe$, la empresa incurre en pérdidas. Para mantener la tarifa de primer mejor se requiere un subsidio o un cargo fijo que cubra la diferencia.}

\sol{78}{B}{El subsidio cruzado combina un objetivo de equidad distributiva con una lógica de precios Ramsey: los usuarios de mayor consumo o menor elasticidad pagan un margen mayor. Puede distorsionar el consumo.}

\sol{79}{D}{El subsidio implícito por m$^3$ es $(costo-precio)$; para 10 m$^3$ es $(costo-precio)\times10$.}

\sol{80}{A}{Subsidio $=(3-1.50)\times8=1.50\times8=12$ soles.}

\sol{81}{C}{Si 30\% del agua producida no se factura, se factura 70\%. El costo por m$^3$ facturado es $2/0.70=2.857\approx2.86$.}

\sol{82}{B}{Con ANF de 20\%, el costo por m$^3$ facturado es $2/0.80=2.50$. La reducción es $2.857-2.50=0.357\approx0.36$.}

\sol{83}{D}{Ratio $=60/1500=0.04=4\%$.}

\sol{84}{A}{Excedente $=(5-2)\times10=30$ soles.}

\sol{85}{C}{La CV es la cantidad de ingreso que, a los nuevos precios, devuelve al consumidor a su utilidad original. La EV se mide a los precios originales.}

\sol{86}{B}{Si la demanda es perfectamente inelástica, la cantidad no cambia; el mayor precio eleva el gasto de los usuarios y el ingreso tarifario.}

\sol{87}{D}{El precio tope tipo RPI-X fija un límite de crecimiento de tarifas; la empresa retiene las ganancias de eficiencia durante el período regulatorio, lo que incentiva reducción de costos.}

\sol{88}{A}{Con la restricción de autofinanciamiento, la solución de segundo mejor es Ramsey-Boiteaux: márgenes inversamente relacionados con las elasticidades.}

\sol{89}{C}{$\Delta\ln Q\approx b\Delta\ln P+c\Delta\ln Y=(-0.4)(10\%)+(0.2)(5\%)=-4\%+1\%=-3\%$.}

\sol{90}{B}{El bloque de consumo básico no identifica perfectamente la pobreza: hogares no pobres con bajo consumo pueden recibir el subsidio (filtración) y hogares pobres con consumo superior al bloque pueden quedar excluidos (subcobertura).}
